% =============================================================================================
% MASTER SCREENPLAY & MATHEMATICAL FOUNDATION
% Highly Accurate Protein Structure Prediction with AlphaFold (AlphaFold 2)
% Seminar: Machine Learning for Genomics (263-5351-00L), ETH Zürich
% Presenters: Evita Bosch Pokunevich (Part I) & Ben Bullinger (Part II)
% =============================================================================================

\documentclass[10pt,a4paper]{article}
\usepackage[margin=20mm,bottom=24mm,footskip=12mm]{geometry}
\usepackage[T1]{fontenc}
\usepackage{XCharter}
\usepackage[xcharter,vvarbb]{newtxmath}
\usepackage{amsmath,bm,booktabs,microtype,xcolor,graphicx,tikz,array}
\usetikzlibrary{arrows.meta,calc,positioning}
\usepackage[colorlinks,linkcolor=ETHBlue,urlcolor=ETHBlue,citecolor=ETHRed]{hyperref}

% --- ETH COLOR PALETTE ---
\definecolor{ETHBlue}{RGB}{33,92,175}
\definecolor{ETHRed}{RGB}{183,53,45}
\definecolor{ETHGray}{RGB}{111,111,111}
\definecolor{ETHGreen}{RGB}{44,162,95}
\definecolor{Gold}{RGB}{204,166,61}

% --- PRESENTATION MACROS ---
\newcommand{\SE}{\mathrm{SE}(3)}
\newcommand{\SO}{\mathrm{SO}(3)}
\newcommand{\R}{\mathbb{R}}
\newcommand{\vx}{\vec{\bm{x}}}
\newcommand{\vt}{\vec{\bm{t}}}
\newcommand{\vq}{\vec{\bm{q}}}
\newcommand{\vk}{\vec{\bm{k}}}
\newcommand{\vval}{\vec{\bm{v}}}
\newcommand{\Tg}{T_{\text{global}}}

% --- SCREENPLAY FORMATTING MACROS ---
\newcommand{\evita}[1]{\par\vspace{2mm}\noindent\textbf{\color{ETHBlue}EVITA:}\enspace #1\par}
\newcommand{\ben}[1]{\par\vspace{2mm}\noindent\textbf{\color{ETHRed}BEN:}\enspace #1\par}
\newcommand{\joint}[1]{\par\vspace{2mm}\noindent\textbf{\color{ETHGray}EVITA \& BEN:}\enspace #1\par}

\newcommand{\slideheader}[4]{%
  \par\vspace{6mm}
  \noindent\tikz{%
    \node[fill=ETHBlue!8, draw=ETHBlue!40, line width=0.8pt, rounded corners=3pt, inner sep=6pt, text width=\linewidth-14pt] {
      {\bfseries\large\textcolor{ETHBlue}{Slide #1: #2}} \hfill {\small\bfseries\color{ETHGray}#3 \quad\textbullet\quad #4}
    };
  }\par\vspace{2mm}
}

\newcommand{\shadowbox}[1]{%
  \par\vspace{1.5mm}
  \noindent\tikz{%
    \node[fill=black!3, draw=black!25, line width=0.5pt, rounded corners=2pt, inner sep=5pt, text width=\linewidth-12pt] {
      \footnotesize\textbf{\color{black!75}SLIDE CONTENT:}\enspace #1
    };
  }\par\vspace{2mm}
}

\newcommand{\mathnote}[1]{%
  \par\vspace{1.5mm}
  \noindent\tikz{%
    \node[fill=ETHRed!4, draw=ETHRed!30, line width=0.6pt, rounded corners=2pt, inner sep=5pt, text width=\linewidth-12pt] {
      \footnotesize\textbf{\color{ETHRed}MATHEMATICAL CORE \& DEEP LEARNING PRINCIPLES:}\par\vspace{1mm}
      \small #1
    };
  }\par\vspace{2mm}
}

\newcommand{\technicalnote}[2]{%
  \par\vspace{1.5mm}
  \noindent\tikz{%
    \node[fill=ETHBlue!4, draw=ETHBlue!30, line width=0.6pt, rounded corners=2pt, inner sep=5pt, text width=\linewidth-12pt] {
      \footnotesize\textbf{\color{ETHBlue}#1:}\par\vspace{1mm}
      \small #2
    };
  }\par\vspace{2mm}
}

\title{\vspace{-8mm}\bfseries\LARGE AlphaFold 2: Architecture, Geometry, and Mathematics\\[2mm]
\large Mathematical Foundations and Presentation Companion Notes}
\author{\textbf{Evita Bosch Pokunevich} \quad\&\quad \textbf{Ben Bullinger}\\[1mm]
\small Machine Learning for Genomics (263-5351-00L), ETH Z\"urich \quad\textbullet\quad Prof.~Dr.~Valentina Boeva}
\date{Fall Semester 2026}

\begin{document}
\maketitle

\vspace{-2mm}
\begin{abstract}
\noindent This document provides the mathematical foundations, architectural derivations, and presentation notes for Jumper et al., \textit{Nature} 596 (2021) (\textit{AlphaFold 2}). It details the dual-track Evoformer architecture, $\mathrm{SE}(3)$-invariant geometry in the Structure Module, Frame Aligned Point Error (FAPE), and the physical and evolutionary constraints governing protein structure prediction.
\end{abstract}

\tableofcontents
\vspace{4mm}
\hrule
\vspace{4mm}

% =============================================================================================
\section{Presentation Structure and Timing Budget}
% =============================================================================================

\begin{center}
\footnotesize
\setlength{\tabcolsep}{3.5pt}
\begin{tabular}{@{}r >{\raggedright\arraybackslash}p{4.5cm} l >{\raggedright\arraybackslash}p{5.9cm} r r@{}}
\toprule
Slide & Title & Speaker & Key Concept / Focus & Slide Time & Cumulative \\
\midrule
\multicolumn{6}{@{}l}{\textbf{Part I: Problem Formulation, CASP Benchmark, and the Evoformer Trunk (Evita Bosch Pokunevich)}} \\
Title & Title Slide & Evita & Context, PDB 6Y4F rendering, paper citation & 0:30 & 0:30 \\
1 & The problem: sequence to structure & Evita & Levinthal's paradox, heavy atoms, 50-yr challenge & 1:15 & 1:45 \\
2 & The benchmark: CASP14 & Evita & Blind test, 0.96\,\AA\ median error vs 1.4\,\AA\ bond length & 1:15 & 3:00 \\
3 & The pipeline: MSA in, coordinates out & Evita & End-to-end macro-pipeline, recycling loop & 1:15 & 4:15 \\
4 & Four facts, four design choices & Evita & The architectural leitmotif of the entire talk & 1:30 & 5:45 \\
5 & The two representations & Evita & MSA tensor $\mathbf{m}_{si}$ and Pair matrix $\mathbf{z}_{ij}$ & 1:15 & 7:00 \\
6 & The input: MSA representation & Evita & Evolutionary covariance, co-mutations, epistasis & 1:15 & 8:15 \\
7 & The input: Pair representation & Evita & Graph edges, spatial distance prior, relative pos & 1:15 & 9:30 \\
8 & The Evoformer block & Evita & Dual-track architecture, 48 blocks, axial attention & 1:15 & 10:45 \\
9 & Evoformer: triangle updates & Evita & Triangle inequality, outer product mean & 1:15 & 12:00 \\
10 & Structure module: residues as rigid bodies & Evita & $\mathrm{SE}(3)$ residue frames, residue gas & 1:15 & 13:15 \\
\midrule
\multicolumn{6}{@{}l}{\textbf{Part II: Geometric Invariance, Structure Module, Training, and Frontiers (Ben Bullinger)}} \\
11 & Invariant Point Attention (IPA) & Ben & $\mathrm{SE}(3)$ invariance proof, local coordinate frames & 1:30 & 14:45 \\
12 & Frame Aligned Point Error (FAPE) & Ben & Loss formulation, chirality sensitivity, metric proof & 1:30 & 16:15 \\
13 & Training data: PDB and self-distillation & Ben & 350k synthetic self-distillation, 15\% BERT masking & 1:00 & 17:15 \\
14 & Ablations: the cost of each component & Ben & CASP14 ablations, $-16.3$ non-linear collapse & 1:15 & 18:30 \\
15 & Intermediate structures in the trunk & Ben & Progressive folding, bimodal dynamics (LmrP, ORF8) & 1:15 & 19:45 \\
16 & Confidence: pLDDT against measured lDDT & Ben & Well-calibrated uncertainty ($r=0.76$), IDR detection & 1:15 & 21:00 \\
17 & Limitations: MSA depth and complexes & Ben & $N_{\text{eff}} < 30$ cliff, multi-chain/conformational gap & 1:15 & 22:15 \\
18 & After AlphaFold 2: 2021 to 2025 & Ben & AFDB (214M), Multimer, ESMFold, Nobel Prize 2024 & 1:00 & 23:15 \\
19 & AlphaFold 3: which choices survived & Ben & Generative shift: pair graph kept, diffusion for IPA & 1:15 & 24:30 \\
20 & Summary: four design choices, measured & Ben & Quantitative ablation scorecard, final synthesis & 0:45 & 25:15 \\
\midrule
\multicolumn{6}{@{}l}{\textbf{Part III: Seminar Discussion and References (Shared)}} \\
21 & Exam Question & Evita/Ben & Course examination analysis & 2:00 & 27:15 \\
22 & Questions for the audience & Evita/Ben & Moderated interactive peer debate & 2:30 & 29:45 \\
23--24 & References & Evita/Ben & Primary literature & 1:00 & 30:45 \\
\midrule
\multicolumn{6}{@{}l}{\textbf{Backup: Top 5 Technical Defense Concepts (Highest Expected Return)}} \\
B1--B5 & Backup Slides (1/5 to 5/5) & Ben & IPA proof, FAPE metric, Triangles, Ablations, AF3 & Q\&A & Q\&A \\
\bottomrule
\end{tabular}
\end{center}

\newpage
% =============================================================================================
\section{Part I: Problem Formulation, CASP Benchmark, and the Evoformer Trunk}
% =============================================================================================

\slideheader{Title}{Highly Accurate Protein Structure Prediction with AlphaFold}{Evita}{0:00--0:30}
\shadowbox{Minimalist title slide: paper citation (Jumper et al., \textit{Nature} 2021), authors, ETH institutional affiliation, presenter names and repository link.}

\evita{Good morning everyone, Prof.~Boeva, and fellow researchers. Today, Ben and I present one of the true milestones of contemporary computational biology: \textit{``Highly accurate protein structure prediction with AlphaFold''}, published in Nature in 2021 by John Jumper and the team at DeepMind. When we say AlphaFold today, we mean AlphaFold~2. Our goal is to unpack the rigorous mathematical and physical inductive biases that enabled deep learning to solve a 50-year-old grand challenge in molecular biology.}

% ---------------------------------------------------------------------------------------------
\slideheader{1}{The problem: sequence to structure}{Evita}{0:30--1:45}
\shadowbox{Left: amino acid string `SIFSYITESTGTPS\dots', arrows pointing right, bullet points on experimental timescale vs sequence database explosion. Right: CASP14 Target T1049 (PDB 6Y4F), blue prediction overlaid on green experiment.}

\evita{Let us begin from first principles. A protein begins life as a one-dimensional polymer: an unbranched sequence of amino acids selected from an alphabet of 20 canonical residues. Every residue shares an identical chemical backbone---Nitrogen, Carbon-alpha, and Carbon---and is distinguished only by its sidechain. 

In the thermodynamic hypothesis formulated by Christian Anfinsen in 1973, native protein structures reside at the global free energy minimum. Yet according to Levinthal's paradox, an unbiased random conformational search of a 100-residue chain across its torsional angles would require $10^{77}$ years---longer than the age of the universe. In nature, proteins fold spontaneously in milliseconds.

For structural biologists, determining one structure experimentally via X-ray crystallography or cryo-EM requires months to years of arduous wet-lab labor. As high-throughput sequencing exploded, the UniProt database grew past hundreds of millions of known sequences, while the Protein Data Bank (PDB) remained constrained to around 180,000 resolved structures. The challenge: can machine learning map directly from the 1D primary sequence to the 3D atomic coordinates of all heavy atoms?}

\mathnote{
Let $\mathcal{A}$ denote the alphabet of 20 amino acids. The primary sequence is $\bm{s} = (s_1, \dots, s_N) \in \mathcal{A}^N$. The target output is the set of Cartesian coordinates for all $M$ heavy (non-hydrogen) atoms:
\begin{equation}
\mathbf{X} = \{ \bm{x}_a \in \R^3 \mid a \in \{1, \dots, M\} \}.
\end{equation}
Each residue $i \in \{1, \dots, N\}$ has backbone atoms $(\mathrm{N}_i, \mathrm{C}\alpha_i, \mathrm{C}_i, \mathrm{O}_i)$ and up to four sidechain dihedral angles $\chi_{i,1}, \chi_{i,2}, \chi_{i,3}, \chi_{i,4}$.
}

% ---------------------------------------------------------------------------------------------
\slideheader{2}{The benchmark: CASP14}{Evita}{1:45--3:00}
\shadowbox{CASP14 bar plot: AlphaFold in blue (0.96\,\AA), 14 next best groups in gray (2.8\,\AA). Dashed red line at 1.4\,\AA\ marked `Atomic bond length $\approx 1.4$\,\AA\ (C--C)'. Horizontal y-axis label atop ticks: `Median backbone error (\AA)'.}

\evita{To evaluate whether computational models truly generalize rather than overfit to known structures, the community established CASP: the Critical Assessment of protein Structure Prediction. In CASP, targets are blind: experimental structures solved by crystallography or cryo-EM are withheld from the public until predictions are formally submitted.

Here are the official results from CASP14 in summer 2020. The metric is root-mean-square deviation of backbone $C\alpha$ atoms on the 95\% best-fitting residues. The gray bars show the fourteen next-best computational teams in the world, hovering with median backbone errors around 2.8\,\AA. 

The blue bar is AlphaFold: its median error is 0.96\,\AA. Notice the horizontal dashed line at 1.4\,\AA. That is the length of a single carbon-carbon covalent bond. AlphaFold crossed the fundamental physical threshold of atomic resolution without human intervention.}

\technicalnote{Benchmark Metric Note}{%
\textit{Why is $\mathrm{C}\alpha\text{-RMSD}_{95}$ evaluated rather than raw all-atom RMSD?}\par\vspace{1mm}
All-atom global RMSD is notoriously sensitive to a single flexible terminal tail swinging in solvent: a 10\,\AA\ displacement of a 5-residue tail inflates the global RMSD of an otherwise atomically accurate 500-residue domain. $\text{RMSD}_{95}$ measures the rigid structural core by trimming the 5\% worst outlier residues, reflecting the true coordinate accuracy of the folded fold.}

% ---------------------------------------------------------------------------------------------
\slideheader{3}{The pipeline: MSA in, coordinates out}{Evita}{3:00--4:15}
\shadowbox{Macro architecture schematic: MSA \& Templates $\to$ 48 Evoformer blocks $\to$ 8 Structure Module iterations $\to$ All-atom relaxation. Recycling loop shown with dashed arrow cycling 4 times.}

\evita{How does AlphaFold accomplish this? Look at the overall macro-architecture. AlphaFold abandons legacy two-stage pipelines that first predicted 2D contact distance maps and then handed them off to separate physical simulation software like Rosetta. 

AlphaFold is end-to-end differentiable. The network takes the query sequence, searches genomic databases to assemble a Multiple Sequence Alignment (MSA), retrieves structural templates, and constructs two fundamental data structures: an MSA representation and a Pair representation.

These two arrays pass through 48 blocks of the Evoformer trunk, which perform deep evolutionary and spatial reasoning. The resulting representations are handed to the Structure Module, an 8-block iterative module that directly outputs 3D rigid bodies. Crucially, the entire system recycles its predictions: the output pair representations and coordinates are fed back into the input for 4 complete passes.}

% ---------------------------------------------------------------------------------------------
\slideheader{4}{Four facts, four design choices}{Evita}{4:15--5:45}
\shadowbox{Clean single-line table matching the 4 biological facts to the 4 architectural design choices in AlphaFold 2. Bottom takeaway in blue alert.}

\evita{Here is the central conceptual thesis of our talk. AlphaFold's breakthrough does not stem from throwing raw, unconstrained compute at generic transformers. It succeeds because every structural inductive bias is a direct answer to a concrete physical or biological fact about proteins.

Consider these four facts:
\begin{enumerate}\setlength{\itemsep}{1mm}
  \item \textbf{Fact 1}: An MSA is an unordered set of evolutionarily related sequences, but each sequence is a 1D chain. \textit{Design choice}: Axial attention alternating along rows and columns.
  \item \textbf{Fact 2}: Real 3D distances must close triangles: $d_{ik} \le d_{ij} + d_{jk}$. \textit{Design choice}: Triangle updates on the residue pair graph edges.
  \item \textbf{Fact 3}: A physical molecule is invariant to global rotations and translations: the Lie group $\SE$. \textit{Design choice}: Local residue reference frames and Invariant Point Attention.
  \item \textbf{Fact 4}: Amino acids are stereochemically chiral; mirror reflections are toxic unphysical molecules. \textit{Design choice}: Frame Aligned Point Error (FAPE) loss in local coordinates.
\end{enumerate}
In the second half of our talk, Ben will show that removing any one of these four design choices causes accuracy to degrade severely.}

% ---------------------------------------------------------------------------------------------
\slideheader{5}{The two representations}{Evita}{5:45--7:00}
\shadowbox{Left: MSA grid (sequences by residues, 256 channels). Right: Pair matrix (residues by residues, 128 channels). Unrotated horizontal labels: `$N_{\text{seq}}$ sequences', `$N_{\text{res}}$ residues'. Takeaway: Dual-track representation preserves raw evolutionary signals.}

\evita{On the pipeline diagram, we saw that the Evoformer operates on two parallel data streams. Before we examine how the network processes them, let us look inside the network's working memory to understand what these two representations actually encode.

First, the MSA representation, $\mathbf{m}_{si} \in \R^{256}$. Here, row index $s$ indexes the sequence in the alignment, column index $i$ indexes the residue position, and each cell holds a 256-dimensional feature vector. The first row ($s=1$) is the query sequence.

Second, the Pair representation, $\mathbf{z}_{ij} \in \R^{128}$. Here, row $i$ and column $j$ index pairs of residues along the protein chain. Each entry $\mathbf{z}_{ij}$ represents the geometric and physical relationship between residue $i$ and residue $j$---such as expected spatial distance, chemical compatibility, or orientation.

The fundamental architectural principle here is the dual-track design. Earlier systems collapsed the multiple sequence alignment into static pairwise contact tables before running neural networks, permanently discarding individual sequence details. AlphaFold maintains both representations simultaneously throughout its entire trunk so that raw evolutionary signals are never lost.}

% ---------------------------------------------------------------------------------------------
\slideheader{6}{The input: MSA representation}{Evita}{7:00--8:15}
\shadowbox{Left: Alignment matrix with correlated mutations in blue/red (salt-bridge charge conservation). Right: Contact extraction from evolutionary variation; raw alignment vs legacy DCA covariance tables.}

\evita{Why does the MSA contain 3D spatial information? The answer lies in evolutionary co-variation. When a protein folds in a living cell, residues that are distant along the primary 1D chain may sit in direct physical contact in 3D space to stabilize a hydrophobic core or form a salt bridge.

If residue $i$ mutates from a small hydrophobe to a bulky residue, the adjacent residue $j$ in 3D space experiences intense selective pressure to mutate compensatorily to avoid steric clash. In 2011, Marks, Colwell, and colleagues demonstrated that analyzing correlated mutations across evolutionary history can reconstruct protein folds.

Prior models, including AlphaFold~1, compressed the MSA down to summary statistics: a 1D position-specific scoring matrix and 2D covariance tables. AlphaFold~2 retains the raw aligned sequences. This allows deep attention mechanisms to track high-order epistasis that summary statistics discard.}

% ---------------------------------------------------------------------------------------------
\slideheader{7}{The input: Pair representation}{Evita}{8:15--9:30}
\shadowbox{Left: Complete residue graph edge features $\mathbf{z}_{ij} \in \mathbb{R}^{128}$ (Fig.~3b). Right: Relative sequence separation and homologous template distance initialization.}

\evita{Now consider the pair representation. Mathematically, $\mathbf{z}$ represents the edge features of a complete graph whose nodes are the residues.

At initialization, what do we know about residue pairs before seeing 3D structure? We know their primary sequence separation: $|i - j|$. We know whether residue $i$ precedes residue $j$. AlphaFold encodes this relative sequence distance through a one-hot clipped positional embedding. Furthermore, if homologous experimental structures exist in the PDB, their $C\beta$--$C\beta$ distance maps and backbone dihedral angles are projected into the pair channels as template information.

Throughout the network, the pair representation serves as a spatial blackboard where geometric constraints are continuously refined.}

% ---------------------------------------------------------------------------------------------
\slideheader{8}{The Evoformer block}{Evita}{9:30--10:45}
\shadowbox{Detailed schematic of one Evoformer block: Row-wise attention with pair bias $\to$ Column-wise attention $\to$ Transition MLP $\to$ Outer product mean $\to$ Triangular multiplicative updates $\to$ Triangular self-attention $\to$ Transition MLP.}

\evita{Now that we have examined both representations, how does the network actually process and refine them? This is the core engine: the Evoformer block. AlphaFold stacks 48 identical Evoformer blocks in series.

Notice the dual-track structure. Information flows simultaneously through the MSA track at the top and the Pair track at the bottom. But they do not operate in isolation:
First, the pair track injects geometric bias into the MSA attention logits.
Second, the MSA track aggregates sequence covariation and writes directly back into the pair edges via an Outer Product Mean.

Let us examine the exact mathematical operations executing inside these blocks.}

\slideheader{9}{Evoformer: triangle updates}{Evita}{10:45--12:00}
\shadowbox{Residue triangle graph ($i, j, k$). Outer product mean injects covariance. Triangle multiplicative updates marginalize paths $\sum_k \mathbf{z}_{ik} \odot \mathbf{z}_{jk}$.}

\evita{Within the pair representation, AlphaFold enforces 3D geometric consistency directly in 2D space. A matrix of pairwise distances in 3D Euclidean space cannot have arbitrary independent entries; it must strictly obey the metric triangle inequality: $d_{ij} \le d_{ik} + d_{jk}$. Standard 2D convolutions or axial self-attention fail to enforce metric closure across triplets.

AlphaFold enforces this graph-geometric consistency by updating edge $(i,j)$ by marginalizing two-hop paths through all intermediate residues $k$:
\begin{equation}
\mathbf{z}_{ij} \leftarrow \mathbf{z}_{ij} + \sum_{k=1}^{N_{\text{res}}} \mathbf{a}_{ik} \odot \mathbf{b}_{jk}.
\end{equation}
Two edges of a triangle combine via Hadamard product to update the third edge. Triangle self-attention further updates edge $(i,j)$ with attention biased by $(j,k)$. Triangle updates enforce metric spatial constraints before 3D coordinates exist.}

% ---------------------------------------------------------------------------------------------
\slideheader{10}{Structure module: residues as rigid bodies}{Evita}{12:00--13:15}
\shadowbox{Residues parameterized by $\mathrm{SE}(3)$ rigid bodies $T_i = (R_i, \vec{\mathbf{t}}_i)$. Sidechain torsion angles $\chi$. Residue gas unconstrained representation.}

\evita{We now transition to coordinate generation: the Structure Module. Instead of treating atoms as arbitrary point clouds or building the protein sequentially from N- to C-terminus, AlphaFold represents each residue as an independent rigid-body frame:
\begin{equation}
T_i = (R_i, \vec{\mathbf{t}}_i) \in \SE,
\end{equation}
where $\vec{\mathbf{t}}_i \in \R^3$ is the position of the Carbon-alpha atom, and rotation matrix $R_i \in \SO$ defines the orientation of the peptide backbone plane. Up to four sidechain dihedral angles $\chi$ place sidechain heavy atoms.

All residues update simultaneously in local coordinates as an unconstrained ``residue gas''. Covalent backbone bond lengths and peptide stereochemistry are not rigidly hardcoded into kinematics; they are learned through end-to-end loss gradients.}

\newpage
% =============================================================================================
\section{Part II: Geometric Invariance, Structure Module, Training, and Frontiers}
% =============================================================================================

\slideheader{11}{Invariant Point Attention (IPA)}{Ben}{13:30--15:00}
\shadowbox{Left: Isometric 3D triad of residue frames $i$ and $j$ with 4 projected points each. Right: Attention formula with Euclidean distance squared in exponent, local coordinate transformation $T_i^{-1} \circ T_j$. Bottom: Takeaway in blue alert.}

\ben{We now arrive at the generative heart of AlphaFold: the Structure Module. 

To predict 3D coordinates, machine learning faces a fundamental geometric requirement: physical molecules are invariant to rigid motions. If you rotate or translate an entire protein in 3D space, its internal free energy and biological identity are strictly unchanged: the symmetry group is the special Euclidean group, $\SE = \SO \ltimes \R^3$. 

AlphaFold does not represent atoms as arbitrary point clouds. It models every residue along the backbone as an independent rigid body frame:
\begin{equation}
T_i = (R_i, \vec{\mathbf{t}}_i) \in \SE,
\end{equation}
where $\vec{\mathbf{t}}_i \in \R^3$ is the position of the Carbon-alpha atom, and rotation matrix $R_i \in \SO$ defines the local orientation of the peptide plane.

Inside the Structure Module, residues communicate via Invariant Point Attention (IPA). Standard attention uses scalar queries and keys. IPA supplements scalar attention with 3D points placed in local residue coordinates:
\begin{equation}
\vec{\mathbf{q}}_{i}^{h,p} \in \R^3, \quad \vec{\mathbf{k}}_{j}^{h,p} \in \R^3.
\end{equation}
These points are transformed into global space by the residue's current pose: $\vec{\mathbf{x}}_i = T_i \circ \vec{\mathbf{q}}_i = R_i \vec{\mathbf{q}}_i + \vec{\mathbf{t}}_i$. The attention weight depends directly on their Euclidean distance in 3D space:
\begin{equation}
a_{ij} = \mathrm{softmax}_j\left( \frac{1}{\sqrt{c}} \mathbf{q}_i^\top \mathbf{k}_j + b_{ij} - \frac{\gamma}{2} \sum_{p=1}^{P} \left\| T_i \circ \vec{\mathbf{q}}_i^{p} - T_j \circ \vec{\mathbf{k}}_j^{p} \right\|^2 \right).
\end{equation}
Notice the mathematical invariance: if a global rotation and translation $\Tg \in \SE$ is applied to all frames, the global distance satisfies:
\begin{equation}
\| (\Tg \circ T_i) \circ \vec{\mathbf{q}}_i - (\Tg \circ T_j) \circ \vec{\mathbf{k}}_j \| = \| \Tg \circ (T_i \circ \vec{\mathbf{q}}_i - T_j \circ \vec{\mathbf{k}}_j) \| = \| T_i \circ \vec{\mathbf{q}}_i - T_j \circ \vec{\mathbf{k}}_j \|.
\end{equation}
The attention weights are rigorously invariant. When value points are aggregated, they are pulled back into frame $i$ via $T_i^{-1} \circ T_j$, where $\Tg$ meets its inverse and cancels exactly. The invariance is mathematically exact; it does not need to be learned.}

\mathnote{
\textbf{The $\SE$ Group Composition Law}:
Let $T = (R, \bm{t}) \in \SE$. The action on a point $\bm{x} \in \R^3$ is $T \circ \bm{x} = R\bm{x} + \bm{t}$. 
The composition of two frames $T_1 \circ T_2$ is:
\begin{equation}
(R_1, \bm{t}_1) \circ (R_2, \bm{t}_2) = (R_1 R_2, R_1 \bm{t}_2 + \bm{t}_1).
\end{equation}
The inverse is $T^{-1} = (R^\top, -R^\top \bm{t})$. 
In the update step, $\Delta T_i = (\Delta R_i, \Delta \bm{t}_i)$ is predicted from invariant scalar features $\mathbf{s}_i$ via the exponential map $\Delta R_i = \exp(\hat{\bm{\omega}}_i)$, and composed locally: $T_i^{(k+1)} = T_i^{(k)} \circ \Delta T_i$, guaranteeing strict $\SE$ equivariance.
}

% ---------------------------------------------------------------------------------------------
\slideheader{12}{Frame Aligned Point Error (FAPE)}{Ben}{15:00--16:30}
\shadowbox{Diagram: Frame $k$ in red, predicted atom in blue, true atom in black, dashed red distance line $d_{ki}$. FAPE loss formula. Clamping boundary at 10\,\AA.}

\ben{Once the network outputs 3D coordinates, how do we train it? 

One might naively propose standard RMSD. But RMSD is fatal for training neural networks: computing RMSD requires performing a global Kabsch alignment (via singular value decomposition) at every training step. SVD alignment is numerically fragile, computationally slow, and worst of all, global alignment gradients allow an error in one distant flexible loop to perturb gradients across the entire stable core.

AlphaFold's solution is FAPE: Frame Aligned Point Error. 
Instead of aligning the prediction to the ground truth once globally, FAPE evaluates every atom from the perspective of \textit{every single residue frame}:
\begin{equation}
\mathcal{L}_{\text{FAPE}} = \frac{1}{N_{\text{frames}} N_{\text{atoms}}} \sum_{k=1}^{N_{\text{frames}}} \sum_{i=1}^{N_{\text{atoms}}} \min\left( d_{\text{clamp}}, \left\| T_k^{-1} \circ \vec{\mathbf{x}}_i - (T_k^{\text{true}})^{-1} \circ \vec{\mathbf{x}}_i^{\text{true}} \right\| \right).
\end{equation}
Notice three profound mathematical properties of FAPE:
\begin{enumerate}\setlength{\itemsep}{1mm}
  \item \textbf{Global Invariance}: Because atoms are transformed into local frame coordinates via $T_k^{-1}$, FAPE is identically invariant to global rigid motions of either the prediction or the ground truth.
  \item \textbf{Clamping Threshold}: The error is clamped at $d_{\text{clamp}} = 10$\,\AA. If two domains are misoriented by 30\,\AA, the gradient saturates to zero, preventing distant domain errors from destroying fine-grained atomic packing within a domain.
  \item \textbf{Chirality Sensitivity}: Standard distance matrices are invariant to mirror reflections. But frames in AlphaFold are constructed via Gram--Schmidt: $\bm{e}_1$ along $\mathrm{C}\alpha \to \mathrm{C}$, $\bm{e}_2$ in the peptide plane, and $\bm{e}_3 = \bm{e}_1 \times \bm{e}_2$. The cross product forces every frame to be right-handed. Under a reflection, the atoms flip but the frames remain right-handed, causing the local $z$-coordinate to negate and producing a massive FAPE penalty. FAPE enforces true L-amino acid stereochemical chirality.
\end{enumerate}}

\technicalnote{Mathematical Foundation Note}{%
\textit{Is FAPE a genuine distance metric?}\par\vspace{1mm}
In Euclidean coordinate space $\R^{3M}$, FAPE is a \textit{pseudometric}: it satisfies non-negativity ($\mathcal{L} \ge 0$), symmetry ($\mathcal{L}(X, Y) = \mathcal{L}(Y, X)$), and the triangle inequality ($\mathcal{L}(X, Z) \le \mathcal{L}(X, Y) + \mathcal{L}(Y, Z)$). However, $\mathcal{L}(X, Y) = 0$ implies $X$ and $Y$ coincide up to a proper rigid transformation in $\SE$. Therefore, on the quotient manifold $\R^{3M} / \SE$, FAPE constitutes a genuine metric.}

% ---------------------------------------------------------------------------------------------
\slideheader{13}{Training data: PDB and self-distillation}{Ben}{16:15--17:15}
\shadowbox{Training dataset composition: 170k PDB structures + 350k synthetic self-distilled Uniclust30 predictions (75\% batch share). 15\% BERT-style masked MSA training.}

\ben{How did AlphaFold train a network with 93 million parameters given that the experimental PDB contains only around 170,000 structures?

DeepMind implemented noisy student self-distillation: a trained model predicted structures for 350,000 unlabelled genomic sequences from Uniclust30. High-confidence predictions were merged into training, comprising 75\% of every batch. This tripled the effective training set size and regularized the network.

Simultaneously, a 15\% masked-MSA loss forces the Evoformer to predict masked amino acids from context, training evolutionary representations directly on sequence databases without requiring structural labels.}

% ---------------------------------------------------------------------------------------------
\slideheader{14}{Ablations: the cost of each component}{Ben}{17:15--18:30}
\shadowbox{CASP14 GDT and PDB lDDT ablations (Fig.~4a). Individual losses vs $-16.3$ catastrophic collapse when removing both IPA and recycling simultaneously.}

\ben{The ablation study in Figure~4a quantitatively validates every single structural inductive bias:
Removing sequence attention costs $2.2$ GDT.
Removing triangle updates costs $4.6$ GDT.
Removing IPA costs $2.3$ GDT.
Removing recycling costs $4.5$ GDT.
Without FAPE, reflection-invariant losses produce 50\% inverted chiral mirror images.

Most remarkably, the components exhibit severe non-linear coupling: removing \textit{both IPA and recycling simultaneously} causes a catastrophic $-16.3$ GDT collapse. Without recycling, IPA lacks the iterative depth to untangle complex knotted topologies; without IPA, recycling cannot project invariant coordinates. Structural inductive biases act synergistically.}

% ---------------------------------------------------------------------------------------------
\slideheader{15}{Trunk intermediate structures}{Ben}{18:45--20:00}
\shadowbox{Fig.~4b redrawn: GDT accuracy plotted over 192 Evoformer blocks across 4 runs. Green curve ORF8, Blue/Orange LmrP. Horizontal y-axis label atop axis: `GDT'.}

\ben{What is actually happening inside the Evoformer trunk during inference? In Figure~4b of the paper, DeepMind froze the trunk and attached a readout structure module after every single block across four recycling passes (192 blocks total).

Look at the trajectories:
For target LmrP (in blue), the structure forms almost immediately: by block 25 of the very first recycling pass, the Global Distance Test (GDT) score exceeds 85. The remaining 160 blocks simply polish sidechain packing.

Now look at target ORF8 (in green), an elusive SARS-CoV-2 protein. At block 48, GDT is barely 40---an unrecognizable disordered tangle. It requires nearly 150 blocks of iterative triangle updates before the fold suddenly snaps into place. This demonstrates that AlphaFold is executing non-trivial iterative spatial refinement, not simple table lookup.}

% ---------------------------------------------------------------------------------------------
\slideheader{16}{Confidence: pLDDT against measured lDDT}{Ben}{20:00--21:15}
\shadowbox{Scatter plot: Predicted mean pLDDT vs Measured lDDT-C$\alpha$ ($n=10,795$ chains). Linear fit line (slope 0.997, $r=0.76$). Horizontal y-axis label atop axis: `measured lDDT-C$\alpha$'.}

\ben{In real biological research, knowing \textit{when} to trust a model is as important as the prediction itself. 

AlphaFold outputs a calibrated confidence score per residue called pLDDT: predicted Local Distance Difference Test. The network does not predict a single scalar; a linear head over the final single representation predicts a categorical probability distribution over 50 discrete bins of local $C\alpha$ accuracy.

Look at the scatter plot over 10,795 experimental chains: the Pearson correlation between predicted pLDDT and experimentally measured lDDT is $r = 0.76$, with a linear regression slope of $0.997$---an almost perfect 1-to-1 diagonal calibration. 

When pLDDT $> 90$, sidechain packing is reliable for computational drug design. When pLDDT $< 50$, the region is not necessarily a model failure; in most cases, it corresponds to an intrinsically disordered protein region (IDR) that has no fixed structure in isolation.}

% ---------------------------------------------------------------------------------------------
\slideheader{17}{Limitations: MSA depth and complexes}{Ben}{21:15--22:30}
\shadowbox{Accuracy vs MSA depth curve ($N_{\text{eff}}$). Vertical red dashed line at 30. Steep drop below 30. Horizontal y-axis label atop axis: `lDDT-C$\alpha$'.}

\ben{Where does AlphaFold 2 break down? 

First: Multiple Sequence Alignment depth. Look at the curve of accuracy against $N_{\text{eff}}$, the number of non-redundant effective sequences in the alignment. Above 30 related sequences, accuracy is high and plateaued. But as soon as $N_{\text{eff}}$ drops below 30, accuracy plummets off a cliff. For orphan proteins, de novo designed proteins, or rapidly evolving antibody loops without an evolutionary family, AlphaFold 2 struggles.

Second: AlphaFold 2 was trained exclusively on isolated, static monomeric chains. But proteins in a living cell are dynamic: they undergo allosteric conformational switches, and over 60\% operate as multi-protein complexes. Predicting static monomers cannot capture the full biochemical reality.}

% ---------------------------------------------------------------------------------------------
\slideheader{18}{After AlphaFold 2: 2021 to 2025}{Ben}{22:30--23:30}
\shadowbox{Timeline from 2021 to 2025: RoseTTAFold, 98.5\% Human Proteome (AFDB), AlphaFold-Multimer, ESMFold, AlphaFold 3, Nobel Prize in Chemistry (2024), AlphaGenome (2025).}

\ben{The release of AlphaFold 2 in 2021 triggered an unprecedented revolution across biological science. 

In partnership with EMBL-EBI, DeepMind released the AlphaFold Protein Structure Database, expanding from the human proteome to over 214 million predicted structures---virtually all cataloged proteins known to science. 

Soon after, AlphaFold-Multimer addressed protein-protein complexes. Meta AI released ESMFold, demonstrating that large language models trained on protein sequences could predict structures in seconds without building MSAs. 

In 2024, Demis Hassabis and John Jumper were awarded the Nobel Prize in Chemistry alongside David Baker. And earlier this year, the architecture was generalized to AlphaGenome, predicting non-coding variant regulatory impact directly from DNA.}

% ---------------------------------------------------------------------------------------------
\slideheader{19}{AlphaFold 3: which choices survived}{Ben}{23:30--24:45}
\shadowbox{Table comparing the 4 AF2 choices against AlphaFold 3 (Nature 2024): MSA raw attention dropped; Triangles kept (48 blocks); Frames/IPA dropped for Diffusion; FAPE dropped for Aligned MSE. Chirality violations 4.4\%.}

\ben{In 2024, DeepMind published AlphaFold~3, predicting complexes of proteins, DNA, RNA, and drug ligands. Looking back at our four design choices from Slide 4, which ones survived in AlphaFold~3?

Notice the verdict:
\begin{enumerate}\setlength{\itemsep}{1mm}
  \item \textbf{Choice 1 (MSA attention)}: Largely dropped. AF3 uses a small 4-block MSA module and does not maintain the MSA array, proving that pair representations carry the true geometric load.
  \item \textbf{Choice 2 (Triangle updates)}: Kept completely intact! The 48-block Pairformer retains triangle multiplicative updates and attention. The spatial metric inductive bias remains indispensable.
  \item \textbf{Choice 3 (Frames and IPA)}: Dropped. Replaced by a 3D score-based Diffusion module operating on unconstrained Cartesian coordinates.
  \item \textbf{Choice 4 (FAPE)}: Dropped. Replaced by an aligned mean squared error diffusion loss.
\end{enumerate}
This reflects Rich Sutton's \textit{Bitter Lesson}: generic generative diffusion scales better across diverse chemical modalities than rigid-body protein frames. 

\textbf{However, there is a catch}: by abandoning right-handed local frames and FAPE, AlphaFold~3 lost exact $\SE$ equivariance and stereochemical chirality! AF3 produces 4.4\% chirality violations and relies heavily on multi-seed ranking and physical clash filters to discard unphysical mirror reflections.}

% ---------------------------------------------------------------------------------------------
\slideheader{20}{Summary: four design choices, measured}{Ben}{24:45--25:30}
\shadowbox{Component ablation scorecard: No raw MSA ($-2.2$), No triangles ($-4.6$), No IPA ($-2.3$), No recycling ($-4.5$), No frames (mirrored structures). Final takeaway sentence.}

\ben{To conclude: the four biological facts we introduced at the start became four mathematical design choices. The ablation study proves their necessity:
Removing sequence attention costs $2.2$ GDT. 
Removing triangle updates costs $4.6$ GDT. 
Removing IPA costs $2.3$ GDT. 
Removing recycling costs $4.5$ GDT.
And removing frames causes the model to produce mirror-image molecules.

AlphaFold 2 solved protein structure prediction not by ignoring biology, but by translating the geometry of physics and evolution into end-to-end differentiable neural architecture. Thank you.}

\newpage
% =============================================================================================
\section{Part III: Seminar Discussion and References}
% =============================================================================================

\slideheader{21}{Exam Question}{Evita \& Ben}{25:30--27:30}
\shadowbox{Left: Assigned course examination prompt. Right: Structured analysis framework (inductive bias, loss mechanics, and genomics).}

\joint{In accordance with course seminar requirements, this slide addresses the examination prompt assigned prior to the talk, applying the mathematical and architectural principles established throughout the presentation to evaluate inductive biases and failure modes.}

% ---------------------------------------------------------------------------------------------
\slideheader{22}{Questions for the audience}{Evita \& Ben}{27:30--30:00}
\shadowbox{Three discussion questions: (1) Inductive bias vs unconstrained diffusion, (2) Has the network learned physics or evolution?, (3) pLDDT reliability in genomics. Bottom: repository link.}

\evita{\textbf{Question 1}: \textit{AlphaFold 2 hardcoded $\SE$ equivariance and stereochemical chirality into its architecture and loss. AlphaFold 3 dropped both in favor of unconstrained diffusion. When is geometric inductive bias indispensable, and when is unconstrained generative modeling preferable?}}

\technicalnote{Analytical Perspective}{%
Inductive bias is critical when data is scarce (e.g., PDB scale, $\sim 10^5$ structures). Unconstrained diffusion dominates when scaling across heterogeneous multi-modal complexes (proteins, nucleic acids, small molecules), provided post-hoc filtering or ranking across random seeds is computationally affordable.}

\ben{\textbf{Question 2}: \textit{Given that accuracy drops precipitously below 30 related sequences, has AlphaFold truly learned the physical energy landscape of protein folding, or has it learned a statistical compression of evolutionary variation? What experiment would settle this?}}

\technicalnote{Analytical Perspective}{%
De novo designed proteins (Baker Lab RFdiffusion) fold accurately without natural evolutionary MSAs, suggesting partial physical energy learning. However, true ab initio folding from single sequences remains the domain of protein language models (ESMFold), indicating AF2 relies heavily on co-evolutionary covariance.}

\evita{\textbf{Question 3}: \textit{pLDDT is trained on the model's own errors. On which proteins would you distrust it, and what should a genomics pipeline do with a low-confidence region?}}

\technicalnote{Analytical Perspective}{%
pLDDT can over-predict confidence on isolated monomer models of obligate multimers that require binding partners to stabilize. In clinical genomics pipelines, low-confidence regions ($\text{pLDDT} < 50$) should be flagged as candidate intrinsically disordered regions (IDRs) rather than discarded as misfolded errors.}

% ---------------------------------------------------------------------------------------------
\slideheader{23--24}{References: Primary Literature \& Machine Learning}{Evita \& Ben}{30:00--31:00}
\shadowbox{Two slides listing standardized citations: Jumper 2021, Tunyasuvunakool 2021, Evans 2021, Baek 2021, Abramson 2024, Marks 2011, Lin 2023, Avsec 2025, Devlin 2019, Xie 2020, Sutton 2019.}

\joint{Complete literature provenance and benchmark publications are cataloged here on Slides 23 and 24.}

\newpage
% =============================================================================================
\section{Mathematical Appendix: Key Derivations (Backup Slides 1--5)}
% =============================================================================================

\noindent This appendix provides step-by-step mathematical proofs and algorithmic foundations corresponding to the 5 technical defense backup slides in \texttt{main.tex}.

% ---------------------------------------------------------------------------------------------
\subsection{Backup (1/5): Exact \texorpdfstring{$\mathrm{SE}(3)$}{SE(3)} Group Invariance of IPA}
\textbf{Slide Counterpart}:\enspace \textsl{Backup (1/5): Exact $\mathrm{SE}(3)$ invariance of IPA (SI 1.8.2, Alg.~22)}

\textbf{Mathematical Derivation}:
The only geometric feature entering the attention affinity weights $a_{ij}$ is the Euclidean distance between points projected into global space from local residue frames $i$ and $j$:
\begin{equation}
\left\| (\Tg \circ T_i) \circ \vq_i - (\Tg \circ T_j) \circ \vk_j \right\|^2 = \left\| \Tg \circ (T_i \circ \vq_i - T_j \circ \vk_j) \right\|^2 = \left\| T_i \circ \vq_i - T_j \circ \vk_j \right\|^2.
\end{equation}
Because the Euclidean norm is identically invariant under any global rigid motion $\Tg = (R_g, \vt_g) \in \SE$, the attention affinity weights $a_{ij}$ are completely unchanged. When vector values are aggregated across all residues $j$, they are transformed back into the local frame of residue $i$:
\begin{equation}
(\Tg \circ T_i)^{-1} \circ \sum_j a_{ij} (\Tg \circ T_j) \circ \vval_j = T_i^{-1} \circ \Tg^{-1} \circ \Tg \circ \sum_j a_{ij} T_j \circ \vval_j = T_i^{-1} \circ \sum_j a_{ij} T_j \circ \vval_j.
\end{equation}
$\Tg$ cancels identically against its own inverse $\Tg^{-1}$. Invariance is an exact algebraic group property; it requires no data-driven learning or numerical approximation.

% ---------------------------------------------------------------------------------------------
\subsection{Backup (2/5): Chirality and FAPE Metric Bounds}
\textbf{Slide Counterpart}:\enspace \textsl{Backup (2/5): Chirality and FAPE metric bounds (SI 1.8.1, 1.8.4, 1.9.4)}

\textbf{Mathematical Derivation}:
Residue coordinate frames are constructed via Gram--Schmidt orthogonalization on backbone atoms $(\mathrm{N}, \mathrm{C}\alpha, \mathrm{C})$:
$\bm{e}_1 = \frac{\bm{x}_{\mathrm{C}} - \bm{x}_{\mathrm{C}\alpha}}{\|\bm{x}_{\mathrm{C}} - \bm{x}_{\mathrm{C}\alpha}\|}$, $\bm{e}_2$ in the $\mathrm{N}$--$\mathrm{C}\alpha$--$\mathrm{C}$ plane, and $\bm{e}_3 = \bm{e}_1 \times \bm{e}_2$. 
By definition of the vector cross product, the frame is strictly right-handed. 
If all atomic coordinates undergo spatial inversion (reflection $\bm{x} \mapsto -\bm{x}$), the frames do not reflect: Gram--Schmidt continues to construct right-handed frames. In local frame coordinates, the $z$-coordinate negates: $(T_i^{-1} \circ \bm{x}_j)_z \mapsto -(T_i^{-1} \circ \bm{x}_j)_z$. 
FAPE evaluates:
\begin{equation}
\mathrm{FAPE}(X, X_{\text{mirror}}) = \frac{2}{Z}\sum_{k,i} \min\big(d_{\text{clamp}}, |(T_k^{-1} \circ \bm{x}_i)_z|\big) > 0.
\end{equation}
This term is strictly positive for any non-planar structure, penalizing the mirrored structure with a massive loss. Models trained with reflection-blind dRMSD produce 50\% mirror images; FAPE eliminates them completely.

\textbf{Optimal Kabsch Metric Lower Bound}:
In Supplementary Section 1.9.4, Jumper et al.\ prove that FAPE is bounded below by optimal Kabsch RMSD:
\begin{equation}
\mathcal{L}_{\text{FAPE}}(X, Y) \ge \min_{R \in \SO, \bm{t} \in \R^3} \sqrt{\frac{1}{M}\sum_{i=1}^M \| R\bm{x}_i + \bm{t} - \bm{y}_i \|^2}.
\end{equation}
FAPE reaches zero if and only if predicted and ground truth point sets coincide under a proper rigid motion. Furthermore, FAPE satisfies the triangle inequality on structures, establishing it as a rigorous pseudometric on $\R^{3M}/\SE$.

% ---------------------------------------------------------------------------------------------
\subsection{Backup (3/5): Triangle Updates and 2D Metric Consistency}
\textbf{Slide Counterpart}:\enspace \textsl{Backup (3/5): Triangle updates and 2D metric consistency (SI 1.6.4--1.6.6)}

\textbf{Algorithmic Formulation}:
In 3D Euclidean space, pairwise distances cannot vary independently; they are strictly constrained by the triangle inequality: $d_{ik} \le d_{ij} + d_{jk}$. Standard 2D convolutions or axial self-attention fail to enforce metric closure across triplets.
AlphaFold enforces this graph-geometric consistency by updating edge $(i,j)$ through Hadamard product path marginalization over all intermediate nodes $k$:
\begin{equation}
\mathbf{z}_{ij} \leftarrow \mathbf{z}_{ij} + \mathrm{Linear}\Big(\sum_{k=1}^{N_{\text{res}}} \sigma(\mathbf{a}_{ik}) \odot \sigma(\mathbf{b}_{jk})\Big).
\end{equation}
Two edges of a triangle combine to update the third edge. In parallel, Triangle Self-Attention updates edge $(i,j)$ with multi-head attention over node triplets $(i,k)$ biased by $(j,k)$. This enforces 3D metric embeddability directly within the 2D pair graph prior to coordinate generation.

% ---------------------------------------------------------------------------------------------
\subsection{Backup (4/5): Non-linear Ablation Synergies and Recycling}
\textbf{Slide Counterpart}:\enspace \textsl{Backup (4/5): Non-linear ablation synergies and recycling (SI 1.10 and 1.13)}

\textbf{Ablation Analysis}:
The ablation study in Figure~4a and Supplementary Section 1.13 reveals critical non-linear synergies:
\begin{itemize}\setlength{\itemsep}{1.2mm}
  \item \textbf{No IPA alone} ($-2.3$ GDT): Frames placed via single linear projection; loses iterative coordinate refinement.
  \item \textbf{No recycling alone} ($-4.5$ GDT): Single forward pass; pair representation cannot re-embed intermediate distances.
  \item \textbf{No IPA AND no recycling} ($\bm{-16.3}$ \textbf{GDT}): Catastrophic structural collapse.
  \item \textbf{No raw MSA} ($-2.2$ GDT): Replaces raw MSA with summary DCA tables; discards higher-order epistatic mutational couplings.
  \item \textbf{No triangle updates} ($-4.6$ GDT): Pair graph violates 3D metric embeddability.
\end{itemize}
The components act synergistically: recycling provides the iterative trunk depth required for IPA to untangle complex knotted topologies, while IPA produces the geometrically grounded coordinate frames that recycling feeds back into the pair graph.

% ---------------------------------------------------------------------------------------------
\subsection{Backup (5/5): AlphaFold 3: Diffusion vs Inductive Bias}
\textbf{Slide Counterpart}:\enspace \textsl{Backup (5/5): AlphaFold 3: Diffusion vs Inductive Bias (Nature 2024)}

\textbf{Architectural Analysis}:
In AlphaFold 3 (Abramson et al., \textit{Nature} 630 (2024)), DeepMind transitioned from explicit geometric inductive biases to all-atom generative diffusion:
\begin{itemize}\setlength{\itemsep}{1.2mm}
  \item \textbf{Unified Tokenization}: Replaces protein-specific residue frames with universal chemical tokens representing amino acids, nucleic acids (DNA, RNA), post-translational modifications, and small-molecule ligands.
  \item \textbf{Pair Graph Survival}: The 48-block Pairformer retains triangle multiplicative updates and attention, proving that relational pair graph inductive bias remained indispensable.
  \item \textbf{Diffusion vs IPA}: IPA and FAPE were dropped in favor of a 3D score-based diffusion model operating directly on raw Cartesian atom coordinates. Global $\SE$ invariance is achieved during training via random rigid-body rotations rather than architectural enforcement.
  \item \textbf{The Trade-off}: Abandoning right-handed Gram--Schmidt frames re-introduced chirality violations (4.4\% in AF3) and steric clashes. AF3 compensates by rolling out 25 trajectories across 5 random seeds, ranked by confidence and penalized for stereochemical violations.
\end{itemize}

\end{document}
